%=====================================================================
%  《量化辛几何导引》中文译注版 —— 导言区
%  编译方式：xelatex main.tex（两遍）
%  说明：本文档类为 ctexbook（专著体例）。若教授希望改用 article
%        体例，仅需替换 \documentclass 一行，其余环境定义保持兼容。
%=====================================================================
\usepackage{amsmath}
\usepackage{amssymb}
\usepackage{amsthm}
\usepackage{mathtools}
\usepackage{bm}
\providecommand{\mathscr}[1]{\mathcal{#1}}
\usepackage{graphicx}
\usepackage{xcolor}
\usepackage[most]{tcolorbox}
\usepackage{caption}
\usepackage{enumitem}
\usepackage{booktabs}
\usepackage{array}
\usepackage{multirow}
\usepackage{fancyhdr}
\usepackage[unicode,colorlinks=true,linkcolor=blue!45!black,
            citecolor=blue!45!black,urlcolor=blue!45!black]{hyperref}
\usepackage{bookmark}

%---------------------- 版式 ----------------------
\captionsetup{font=small,labelformat=empty,justification=centering,skip=4pt}
\setlist{itemsep=2pt,topsep=3pt,parsep=0pt}
\allowdisplaybreaks[4]
\linespread{1.28}
\setlength{\emergencystretch}{2.5em}
\setlength{\parindent}{2em}

%---------------------- 常用记号 ----------------------
\newcommand{\R}{\mathbb{R}}
\newcommand{\C}{\mathbb{C}}
\newcommand{\Z}{\mathbb{Z}}
\newcommand{\Q}{\mathbb{Q}}
\newcommand{\N}{\mathbb{N}}
\newcommand{\CP}{\mathbb{CP}}
\newcommand{\RP}{\mathbb{RP}}
\newcommand{\HP}{\mathbb{HP}}
\newcommand{\T}{\mathbb{T}}
\newcommand{\dd}{\mathrm{d}}
\newcommand{\ee}{\mathrm{e}}
\newcommand{\ii}{\mathrm{i}}
\newcommand{\jj}{\mathrm{j}}
\newcommand{\kk}{\mathrm{k}}
\newcommand{\id}{\mathrm{id}}
\newcommand{\Op}{\operatorname}
\newcommand{\Crit}{\operatorname{Crit}}
\newcommand{\ind}{\operatorname{ind}}
\newcommand{\Area}{\operatorname{Area}}
\newcommand{\Vol}{\operatorname{Vol}}
\newcommand{\Symp}{\operatorname{Symp}}
\newcommand{\Diff}{\operatorname{Diff}}
\newcommand{\Ham}{\operatorname{Ham}}
\newcommand{\Aut}{\operatorname{Aut}}
\newcommand{\End}{\operatorname{End}}
\newcommand{\Hom}{\operatorname{Hom}}
\newcommand{\coker}{\operatorname{coker}}
\newcommand{\sgn}{\operatorname{sgn}}
\newcommand{\codim}{\operatorname{codim}}
\newcommand{\Sp}{\operatorname{Sp}}
\newcommand{\vol}{\operatorname{vol}}
\newcommand{\Cont}{\operatorname{Cont}}
\newcommand{\SH}{\operatorname{SH}}
\newcommand{\HF}{\operatorname{HF}}
\newcommand{\CF}{\operatorname{CF}}
\newcommand{\CH}{\operatorname{CH}}
\newcommand{\Ob}{\operatorname{Ob}}
\newcommand{\diag}{\operatorname{diag}}
\newcommand{\Mat}{\operatorname{Mat}}
\newcommand{\Maslov}{\operatorname{Maslov}}
\newcommand{\PSH}{\operatorname{PSH}}

%---------------------- 定理类环境 ----------------------
% 编号一律沿用原书（手工传入），确保与原文及交叉引用完全一致。
\newcommand{\whead}[2]{%
  \par\medskip\noindent\hangindent=2.2em\hangafter=1
  \textbf{#1~#2}\quad}
\newenvironment{wdefn}[1]{\whead{定义}{#1}}{\par\medskip}
\newenvironment{wthm}[1]{\whead{定理}{#1}}{\par\medskip}
\newenvironment{wprop}[1]{\whead{命题}{#1}}{\par\medskip}
\newenvironment{wlem}[1]{\whead{引理}{#1}}{\par\medskip}
\newenvironment{wcor}[1]{\whead{推论}{#1}}{\par\medskip}
\newenvironment{wconj}[1]{\whead{猜想}{#1}}{\par\medskip}
\newenvironment{wexam}[1]{\whead{例}{#1}}{\par\medskip}
\newenvironment{wexer}[1]{\whead{练习}{#1}}{\par\medskip}
\newenvironment{wrem}[1]{\whead{注}{#1}}{\par\medskip}
\newenvironment{wnot}[1]{\whead{记号}{#1}}{\par\medskip}
\newenvironment{wclm}[1]{\whead{断言}{#1}}{\par\medskip}
\newenvironment{wprob}[1]{\whead{问题}{#1}}{\par\medskip}
\newenvironment{wque}[1]{\whead{疑问}{#1}}{\par\medskip}
\newenvironment{wfact}[1]{\whead{事实}{#1}}{\par\medskip}
\newenvironment{wproof}[1][]{%
  \par\medskip\noindent
  \if\relax\detokenize{#1}\relax\textbf{证明}\else\textbf{#1}\fi\quad}%
  {\hfill$\square$\par\medskip}

%---------------------- 本节重点（Key Takeaways） ----------------------
\newtcolorbox{keytakeaways}{
  enhanced, breakable,
  colback=blue!3!white, colframe=blue!55!black,
  boxrule=0.6pt, arc=2pt,
  left=8pt, right=8pt, top=4pt, bottom=4pt,
  title={本节重点（Key Takeaways）},
  fonttitle=\bfseries\small, fontupper=\small,
  coltitle=white, colbacktitle=blue!55!black
}

%---------------------- 章节标题格式 ----------------------
\ctexset{
  chapter={
    format=\centering\heiti\LARGE,
    name={第,章},
    number=\arabic{chapter},
    beforeskip=8pt,
    afterskip=22pt
  },
  section={format=\Large\bfseries, beforeskip=18pt, afterskip=8pt},
  subsection={format=\large\bfseries, beforeskip=12pt, afterskip=5pt},
  subsubsection={format=\normalsize\bfseries, beforeskip=8pt, afterskip=3pt}
}

%---------------------- 页眉页脚 ----------------------
\pagestyle{fancy}
\fancyhf{}
\fancyhead[L]{\small 量化辛几何导引}
\fancyhead[R]{\small\leftmark}
\fancyfoot[C]{\thepage}
\renewcommand{\headrulewidth}{0.4pt}
\fancypagestyle{plain}{\fancyhf{}\fancyfoot[C]{\thepage}\renewcommand{\headrulewidth}{0pt}}
\renewcommand{\chaptermark}[1]{\markboth{\chaptername\ \thechapter\quad #1}{}}
\renewcommand{\sectionmark}[1]{\markright{\thesection\quad #1}}
